\documentclass[12pt, twoside]{article}
\input{/Users/chris/github/course-files/docs/ib/preamble}

\fancyhead[LO]{La Scuola d'Italia / Dr.\ Huson / IB Math 12 HL \\* 14 October 2026}
\fancyhead[RO]{\fbox{\begin{minipage}{6cm}\footnotesize Nickname:\\[0.5cm]\end{minipage}}}
\fancyhead[LE]{\fbox{\begin{minipage}{6cm}\footnotesize Nickname:\\[0.5cm]\end{minipage}}}
\fancyfoot[C]{\thepage}

\begin{document}
\subsubsection*{12.6 Quiz: Complex Numbers}

Answer all five problems. Use the spaces provided for Problems 1 and 2.
Answer Problems 3--5 on lined paper.

\begin{enumerate}[itemsep=0.15cm]

%--- Problem 1: Cartesian form; conjugate; product and quotient; modulus and argument [6 marks] ---
%--- Source: Original (freshly authored 2026-10-14, HL syllabus 1.12, 1.13) ---
\item[1.] [Maximum mark: 6]

Let $z_{1} = 3 + 4i$ and $z_{2} = 1 - 2i$.

\begin{enumerate}[itemsep=0.15cm]
  \item Find $z_{1}z_{2}^{*}$, giving your answer in the form $a + bi$,
  where $a, b \in \mathbb{R}$. \hfill [2]
  \item Find $\dfrac{z_{1}}{z_{2}}$, giving your answer in the form
  $a + bi$, where $a, b \in \mathbb{R}$. \hfill [2]
  \item Find the modulus and the argument of $\dfrac{z_{1}}{z_{2}}$.
  \hfill [2]
\end{enumerate}

\begin{tikzpicture}
  \draw (-0.6,0) rectangle (14.6,15.0);
  \foreach \y in {1,2,...,14} \draw[dotted] (0,\y)--(14.1,\y);
\end{tikzpicture}

\newpage

%--- Problem 2: Euler and polar form; product in polar and Cartesian form; geometric meaning of multiplication [7 marks] ---
%--- Source: Original (freshly authored 2026-10-14, HL syllabus 1.13) ---
\item[2.] [Maximum mark: 7]

Let $z = -\sqrt{3} + i$ and
$w = 3\left(\cos\dfrac{\pi}{4} + i\sin\dfrac{\pi}{4}\right)$.

\begin{enumerate}[itemsep=0.15cm]
  \item Write $z$ in the form $re^{i\theta}$, where $r > 0$ and
  $-\pi < \theta \leq \pi$. \hfill [2]
  \item Find $zw$ in the form $r(\cos\theta + i\sin\theta)$, where $r > 0$
  and $-\pi < \theta \leq \pi$. \hfill [2]
  \item Write $zw$ in the form $a + bi$. \hfill [1]
  \item On an Argand diagram, the point representing $z$ is mapped to the
  point representing $zw$. Describe this transformation
  geometrically. \hfill [2]
\end{enumerate}

\begin{tikzpicture}
  \draw (-0.6,0) rectangle (14.6,15.0);
  \foreach \y in {1,2,...,14} \draw[dotted] (0,\y)--(14.1,\y);
\end{tikzpicture}

\newpage

\emph{Answer Problems 3--5 on lined paper.}

%--- Problem 3: Conjugate root theorem; sum and product of roots of a cubic; Argand triangle [13 marks] ---
%--- Source: Original (freshly authored 2026-10-14, HL syllabus 1.12, 1.14; link 2.12) ---
\item[3.] [Maximum mark: 13]

The equation
\[
z^{3} + az^{2} + bz + 26 = 0,
\]
where $a, b \in \mathbb{R}$, has a root $3 + 2i$.

\begin{enumerate}[itemsep=0.15cm]
  \item Write down a second root of the equation, giving a reason for
  your answer. \hfill [2]
  \item By considering the product of the roots, find the third root.
  \hfill [3]
  \item Find the value of $a$ and the value of $b$. \hfill [4]
\end{enumerate}

The three roots are represented by the points A, B and C on an Argand
diagram, where A represents $3 + 2i$ and C represents the real root.

\begin{enumerate}[itemsep=0.15cm]
  \item[(d)] Show that triangle ABC is isosceles. \hfill [2]
  \item[(e)] Find the area of triangle ABC. \hfill [2]
\end{enumerate}


%--- Problem 4: Proof of de Moivre by induction; multiple-angle identity; exact power [14 marks] ---
%--- Source: Original (freshly authored 2026-10-14, HL syllabus 1.14, 1.15; link 3.10) ---
\item[4.] [Maximum mark: 14]

\begin{enumerate}[itemsep=0.15cm]
  \item Use mathematical induction to prove that
  \[
  (\cos\theta + i\sin\theta)^{n} = \cos n\theta + i\sin n\theta
  \]
  for all $n \in \mathbb{Z}^{+}$. \hfill [6]
  \item By using de Moivre's theorem with $n = 4$, show that \\
  $\cos 4\theta = 8\cos^{4}\theta - 8\cos^{2}\theta + 1$. \hfill [5]
  \item Use de Moivre's theorem to find $\left(1 - \sqrt{3}\,i\right)^{10}$,
  giving your answer in the form $a + b\sqrt{3}\,i$, where
  $a, b \in \mathbb{Z}$. \hfill [3]
\end{enumerate}


%--- Problem 5: Euler form; cube roots (rational exponent); sum of roots; w^(2/3) [13 marks] ---
%--- Source: Original (freshly authored 2026-10-14, HL syllabus 1.13, 1.14; link 2.12) ---
\item[5.] [Maximum mark: 13]

Let $w = -4\sqrt{2} + 4\sqrt{2}\,i$.

\begin{enumerate}[itemsep=0.15cm]
  \item Write $w$ in the form $re^{i\theta}$, where $r > 0$ and
  $-\pi < \theta \leq \pi$. \hfill [2]
  \item Solve the equation $z^{3} = w$, giving each solution in the form
  $re^{i\theta}$, where $r > 0$ and $-\pi < \theta \leq \pi$. \hfill [5]
  \item Write the solution that lies in the first quadrant in the form
  $a + bi$. \hfill [1]
  \item Show that the sum of the three solutions is zero. \hfill [2]
  \item Find the three values of $w^{\frac{2}{3}}$, giving each in the form
  $a + bi$, where $a, b \in \mathbb{R}$. \hfill [3]
\end{enumerate}

\end{enumerate}
\end{document}
